\documentclass[aspectratio=169,11pt]{beamer}
\usepackage{../../../shared/theme/esmad_beamer_theme}
\usepackage{amsmath,amssymb}
\usepackage{booktabs}

\title[Selection on Observables]{Selection on Observables}
\subtitle{Matching, Weighting, and Doubly Robust Estimation}
\date{\today}

\newcommand{\E}{\mathbb{E}}

\begin{document}

\begin{frame}
\titlepage
\end{frame}

\begin{frame}{Learning Goals}
\begin{itemize}
\item state the assumptions required to identify treatment effects from observed covariates;
\item explain the role of the propensity score;
\item compare matching, inverse-probability weighting, and outcome regression;
\item diagnose lack of overlap and covariate imbalance;
\item derive the logic of augmented inverse-probability weighting;
\item distinguish double robustness from robustness to unmeasured confounding.
\end{itemize}
\end{frame}

\begin{frame}{Outline}
\tableofcontents
\end{frame}

\section{Identification from Observed Covariates}

\begin{frame}{Conditional Exchangeability}
Let \(D\in\{0,1\}\) be treatment and \(X\) observed pre-treatment covariates.

A standard identification assumption is
\[
(Y(1),Y(0)) \perp D \mid X.
\]

\begin{block}{Meaning}
After conditioning on \(X\), treatment assignment is independent of the potential outcomes.
\end{block}

\begin{alertblock}{What this requires}
All important common causes of treatment and outcome must be adequately captured by \(X\).
\end{alertblock}
\end{frame}

\begin{frame}{Positivity}
We also need
\[
0<P(D=1\mid X=x)<1
\]
for covariate values relevant to the target population.

If some types of units are always treated or never treated, their counterfactual treatment state is unsupported by the data.

\begin{alertblock}{No overlap, no nonparametric identification}
Model extrapolation cannot create empirical support where none exists.
\end{alertblock}
\end{frame}

\begin{frame}{Consistency}
If unit \(i\) receives treatment \(D_i=d\), the observed outcome must correspond to the potential outcome under that treatment:
\[
Y_i=Y_i(d).
\]

This assumes the treatment is well defined.

\begin{block}{Threat}
If “treatment” bundles heterogeneous versions with different causal effects, the estimand may be ill posed.
\end{block}
\end{frame}

\section{The Propensity Score}

\begin{frame}{Definition}
The propensity score is
\[
e(X)=P(D=1\mid X).
\]

Rosenbaum and Rubin showed that under conditional exchangeability,
\[
(Y(1),Y(0))\perp D\mid e(X).
\]

\begin{block}{Role}
The propensity score compresses a potentially high-dimensional adjustment problem into a scalar balancing score.
\end{block}
\end{frame}

\begin{frame}{What the Propensity Score Is Not}
The propensity score is not:
\begin{itemize}
\item a prediction target whose AUC should be maximized at all costs;
\item evidence that unmeasured confounding is absent;
\item a causal effect estimator by itself;
\item useful if treatment probabilities are near zero or one for large parts of the sample.
\end{itemize}

\begin{alertblock}{Priority}
Estimate treatment assignment well enough to achieve balance and overlap for the target estimand.
\end{alertblock}
\end{frame}

\section{Matching}

\begin{frame}{Nearest-Neighbor Matching}
For a treated unit \(i\), find an untreated unit \(j\) with similar covariates or propensity score.

A simple ATT estimator is
\[
\widehat{ATT}
=
\frac{1}{N_1}
\sum_{i:D_i=1}
\left(
Y_i-Y_{m(i)}
\right).
\]

Here \(m(i)\) denotes the matched control.
\end{frame}

\begin{frame}{Calipers and Common Support}
A caliper prevents poor matches:
\[
|\widehat e(X_i)-\widehat e(X_j)|<c.
\]

\begin{itemize}
\item narrow calipers improve similarity but may discard many units;
\item wide calipers retain more units but accept worse matches;
\item discarded units change the population to which the result applies.
\end{itemize}
\end{frame}

\begin{frame}{Matching Diagnostics}
After matching, inspect:
\begin{itemize}
\item standardized mean differences;
\item variance ratios;
\item propensity-score overlap;
\item distributional balance, not just means;
\item number and characteristics of discarded units.
\end{itemize}

\begin{alertblock}{Balance first}
Do not judge matching quality by whether the treatment-effect estimate “looks reasonable.”
\end{alertblock}
\end{frame}

\section{Inverse-Probability Weighting}

\begin{frame}{ATE Weights}
For the ATE:
\[
w_i=
\frac{D_i}{e(X_i)}
+
\frac{1-D_i}{1-e(X_i)}.
\]

Weighted treated and untreated groups target the same covariate distribution.

\begin{block}{Pseudo-population}
IPW reweights the observed sample to mimic a population in which treatment is independent of measured covariates.
\end{block}
\end{frame}

\begin{frame}{ATT Weights}
For the ATT, one common weighting scheme is:
\[
w_i=
\begin{cases}
1, & D_i=1, \\\
\dfrac{e(X_i)}{1-e(X_i)}, & D_i=0.
\end{cases}
\]

The control group is reweighted to resemble the treated group.

\begin{alertblock}{Estimand matters}
ATE and ATT weights are different because they target different populations.
\end{alertblock}
\end{frame}

\begin{frame}{Extreme Weights}
If \(e(X_i)\approx 0\) for a treated unit,
\[
\frac{1}{e(X_i)}
\]
becomes huge.

Similarly, controls with \(e(X_i)\approx 1\) receive huge weights.

Consequences:
\begin{itemize}
\item high variance;
\item effective sample size collapse;
\item sensitivity to a few observations;
\item unstable estimates.
\end{itemize}
\end{frame}

\begin{frame}{Stabilized Weights}
Stabilized weights can reduce variance:
\[
SW_i=
\frac{P(D=D_i)}
{P(D=D_i\mid X_i)}.
\]

They preserve the basic weighting logic while moderating scale.

\begin{block}{Still required}
Overlap diagnostics and inspection of the full weight distribution remain necessary.
\end{block}
\end{frame}

\section{Balance and Effective Sample Size}

\begin{frame}{Standardized Mean Differences}
For covariate \(X_k\), a standardized mean difference is roughly
\[
SMD_k
=
\frac{\bar X_{k,1}-\bar X_{k,0}}
{S_{pooled}}.
\]

After weighting or matching, good adjustment should substantially reduce imbalance across important covariates.

\begin{alertblock}{No universal magic threshold}
Rules such as 0.1 are useful conventions, not proofs of causal adequacy.
\end{alertblock}
\end{frame}

\begin{frame}{Effective Sample Size}
For weights \(w_i\), a common effective sample-size diagnostic is
\[
ESS=
\frac{(\sum_i w_i)^2}{\sum_i w_i^2}.
\]

When a few units dominate the weighted sample, ESS can be far smaller than the nominal sample size.
\end{frame}

\section{Outcome Regression}

\begin{frame}{Regression Adjustment}
Estimate outcome models
\[
\mu_d(X)=\E[Y\mid D=d,X]
\]
for \(d\in\{0,1\}\).

Then estimate the ATE by
\[
\widehat{ATE}
=
\frac{1}{n}
\sum_i
\left[
\widehat\mu_1(X_i)
-
\widehat\mu_0(X_i)
\right].
\]

\begin{block}{Strength}
Outcome modeling can be efficient when the conditional expectation is modeled well.
\end{block}
\end{frame}

\begin{frame}{Outcome-Model Failure}
A misspecified outcome regression can produce bias even if treatment assignment is fully captured by observed covariates.

Common threats:
\begin{itemize}
\item omitted nonlinearities;
\item missing interactions;
\item extrapolation outside overlap;
\item poor regularization choices.
\end{itemize}
\end{frame}

\section{Doubly Robust Estimation}

\begin{frame}{Augmented IPW}
A common AIPW estimator for the ATE is
\[
\widehat\tau
=
\frac{1}{n}
\sum_i
\Bigg[
\widehat\mu_1(X_i)
-
\widehat\mu_0(X_i)
+
\frac{D_i(Y_i-\widehat\mu_1(X_i))}{\widehat e(X_i)}
-
\frac{(1-D_i)(Y_i-\widehat\mu_0(X_i))}{1-\widehat e(X_i)}
\Bigg].
\]
\end{frame}

\begin{frame}{Why “Doubly Robust”?}
Under standard conditions, the estimator is consistent if either:
\begin{itemize}
\item the propensity model is correctly specified, or
\item the outcome model is correctly specified.
\end{itemize}

\begin{alertblock}{Crucial qualification}
This robustness concerns nuisance-model misspecification. It does not protect against unmeasured confounding, positivity failure, or ill-defined treatment.
\end{alertblock}
\end{frame}

\begin{frame}{Orthogonalization Intuition}
AIPW combines:
\begin{itemize}
\item a model-based prediction of each unit's two potential outcomes;
\item a residual correction weighted by treatment-assignment probabilities.
\end{itemize}

The correction reduces first-order sensitivity to nuisance-estimation errors when regularity conditions hold.
\end{frame}

\section{Trimming and Target Populations}

\begin{frame}{Trimming for Poor Overlap}
One may restrict analysis to
\[
\epsilon
<
\widehat e(X_i)
<
1-\epsilon.
\]

This reduces extreme weights and extrapolation.

\begin{alertblock}{But the estimand changes}
The result now applies to an overlap population, not necessarily the original ATE population.
\end{alertblock}
\end{frame}

\begin{frame}{Overlap Weights}
Overlap weighting emphasizes units with treatment probabilities near 0.5.

This can:
\begin{itemize}
\item improve precision;
\item reduce dominance by extreme observations;
\item target the population with strongest empirical equipoise.
\end{itemize}

Again, the target estimand must be stated explicitly.
\end{frame}

\section{Unmeasured Confounding}

\begin{frame}{The Fundamental Limitation}
All observed-covariate methods rest on the assumption that relevant confounding is measured.

If an unobserved variable \(U\) affects both treatment and outcome,
\[
U\rightarrow D,
\qquad
U\rightarrow Y,
\]
then matching, weighting, and regression on \(X\) need not identify the causal effect.
\end{frame}

\begin{frame}{Sensitivity Analysis}
Sensitivity analysis asks:
\begin{itemize}
\item how strong would an unmeasured confounder need to be to change the conclusion?
\item how much imbalance in an omitted factor would overturn the estimate?
\item how sensitive are results to plausible hidden-bias parameters?
\end{itemize}

\begin{block}{Purpose}
Sensitivity analysis quantifies fragility; it does not remove hidden bias.
\end{block}
\end{frame}

\section{Applied Reasoning}

\begin{frame}{Example: Targeted Marketing Campaign}
Suppose engaged customers are more likely to receive a campaign.

Potential confounders include:
\begin{itemize}
\item prior purchases;
\item site engagement;
\item tenure;
\item demographic characteristics;
\item previous campaign exposure.
\end{itemize}

The causal question is whether campaign exposure changes purchasing, not whether exposed customers purchase more.
\end{frame}

\begin{frame}{A Credible Observational Workflow}
\begin{enumerate}
\item Define treatment, outcome, population, and estimand.
\item Identify pre-treatment confounders from domain knowledge.
\item Estimate the propensity score.
\item Inspect overlap before estimating effects.
\item Match or weight according to the target estimand.
\item Diagnose post-adjustment covariate balance.
\item Inspect weights and effective sample size.
\item Estimate treatment effects with appropriate uncertainty.
\item Compare with outcome regression or AIPW.
\item Run sensitivity analysis for unmeasured confounding.
\end{enumerate}
\end{frame}

\begin{frame}{Common Failure Modes}
\begin{itemize}
\item maximizing propensity-score predictive accuracy instead of covariate balance;
\item matching on post-treatment variables;
\item ignoring lack of overlap;
\item hiding extreme weights;
\item assuming balance proves no hidden confounding;
\item reporting trimmed estimates as if they targeted the original ATE;
\item treating double robustness as immunity to design failure.
\end{itemize}
\end{frame}

\begin{frame}{Synthesis: What Adjustment Can and Cannot Do}
\begin{itemize}
\item Selection-on-observables designs require conditional exchangeability, positivity, and consistency.
\item The propensity score is a balancing device, not a causal effect.
\item Matching and weighting construct more comparable treated and untreated groups.
\item Balance diagnostics evaluate observed covariates after adjustment.
\item Outcome regression models the counterfactual outcome surface directly.
\item AIPW combines weighting and outcome regression with double robustness to nuisance-model misspecification.
\item None of these methods solves important unmeasured confounding without additional structure.
\end{itemize}
\end{frame}

\begin{frame}{Selected References}
\small
\begin{itemize}
\item Rosenbaum, P. R., and Rubin, D. B. (1983). The Central Role of the Propensity Score in Observational Studies for Causal Effects. \textit{Biometrika}, 70(1), 41--55.
\item Hernan, M. A., and Robins, J. M. (2020). \textit{Causal Inference: What If}. Chapman \& Hall/CRC.
\item Imbens, G. W., and Rubin, D. B. (2015). \textit{Causal Inference for Statistics, Social, and Biomedical Sciences}. Cambridge University Press.
\item Austin, P. C. (2011). An Introduction to Propensity Score Methods for Reducing the Effects of Confounding in Observational Studies. \textit{Multivariate Behavioral Research}, 46(3), 399--424.
\item Robins, J. M., Rotnitzky, A., and Zhao, L. P. (1994). Estimation of Regression Coefficients When Some Regressors Are Not Always Observed. \textit{Journal of the American Statistical Association}, 89(427), 846--866.
\end{itemize}
\end{frame}

\contactslide

\end{document}
